\chapter{Troubleshooting}
\label{ch:troubleshooting}

If your project does not work as expected, this chapter helps you identify common problems on
\textbf{both sides of the system}:

\begin{itemize}
    \item \textbf{Android (EmitterApp):} UI controls, Bluetooth pairing/connection, and whether commands/telemetry are being processed.
    \item \textbf{ESP32 (ESP32\_BT\_Controller):} firmware upload, pin wiring, power integrity, Bluetooth SPP link, and telemetry generation.
\end{itemize}

\section{Before You Start: Quick Checklist}
\label{sec:quick-checklist}

\begin{itemize}
    \item Confirm the ESP32 is powered and stable (no brownouts/resets).
    \item Confirm your Android device is paired to the correct Bluetooth name (default is \texttt{ESP32-BT}).
    \item Keep the phone within a few meters for initial testing.
    \item Ensure \textbf{all grounds are common} (ESP32, motor driver, sensors, external 5V/servo supply).
    \item If you are using boot-sensitive pins (e.g., GPIO0), disconnect external circuits temporarily during flashing.
\end{itemize}

\section{ESP32 Firmware Does Not Upload}
\label{sec:esp32-upload}

If firmware does not upload correctly:

\begin{itemize}
    \item Verify the correct board/target is selected in your IDE/tooling.
    \item Confirm the correct serial port is selected.
    \item Try a different USB cable (many ``charge-only'' cables cause upload failures).
    \item Press and hold \textbf{BOOT} during upload if your board requires manual bootloader entry.
    \item Disconnect external wiring from \textbf{boot-strap pins} during flashing (common offenders include GPIO0 and GPIO2).
\end{itemize}

\subsection*{Boot-strap pin notes (important)}
The firmware pin configuration uses GPIO0 as a digital indicator input (\texttt{PIN\_IND4 = GPIO0}) in
\texttt{FULL\_RC\_MODE}. GPIO0 is a boot pin: external pull-ups/pull-downs can prevent boot or flashing.
If flashing fails, unplug anything connected to GPIO0, flash again, then reconnect.

\section{ESP32 Powers On but Resets Randomly}
\label{sec:esp32-resets}

Random resets are usually power-related:

\begin{itemize}
    \item Motors/servos can pull voltage down suddenly. Use a dedicated regulator for the ESP32 and a separate supply for motors/servos.
    \item Ensure the ESP32 supply can provide peak current (Wi-Fi/Bluetooth bursts can cause dips).
    \item Add bulk capacitance near motor drivers/servo rails (e.g., 470--1000~\(\mu\)F) and a smaller ceramic capacitor near the ESP32.
    \item Check for overheating regulators on dev boards when powering servos from the same 5V rail (not recommended).
\end{itemize}

\section{Bluetooth Pairing and Connection Problems (Android \texorpdfstring{$\leftrightarrow$}{<->} ESP32)}
\label{sec:bluetooth-problems}

\subsection*{Symptom: Device not visible in scan}
\begin{itemize}
    \item Confirm the ESP32 is powered and running (LED/activity indicators if available).
    \item Ensure the phone is close to the ESP32.
    \item Power-cycle the ESP32.
\end{itemize}

\subsection*{Symptom: Pairing asks for PIN or fails}
The default PIN is commonly \texttt{1234}.
\begin{itemize}
    \item Remove (``forget'') the device in Android Bluetooth settings, then pair again.
    \item Try pairing in Android system settings first, then connect from the app.
\end{itemize}

\subsection*{Symptom: App says ``Connected'' but controls do nothing}
\begin{itemize}
    \item Verify you connected to the correct device (if multiple \texttt{ESP32-BT} devices exist nearby).
    \item Confirm the ESP32 is running the expected firmware build (not an old sketch).
    \item On the ESP32 side, use serial logging (if enabled) to confirm command packets are arriving.
\end{itemize}

\subsection*{Symptom: Connects, then disconnects repeatedly}
\begin{itemize}
    \item This is often a \textbf{power stability} problem on the ESP32 (see Section~\ref{sec:esp32-resets}).
    \item Reduce motor/servo load during initial testing (disconnect motors, test with LEDs first).
    \item Keep the phone close, and avoid metal enclosures that shield Bluetooth.
\end{itemize}

\section{Controls Feel Wrong in the App (Android Side)}
\label{sec:android-controls}

\subsection*{Joystick returns to center vs stays in place}
EmitterApp supports different joystick modes (for example spring vs hold). If the robot behaves
as if throttle is ``not staying'' or steering is ``stuck'':

\begin{itemize}
    \item Check the selected stick modes in the app settings.
    \item For throttle, many users prefer a \textbf{hold} behavior; for steering, many prefer \textbf{spring}.
\end{itemize}

\subsection*{Buttons/switches appear but do not affect hardware}
\begin{itemize}
    \item Confirm that the ESP32 firmware maps those channels to real pins/outputs in your chosen project mode.
    \item If you use \texttt{FULL\_RC\_MODE\_MCP}, verify MCP23017 wiring and I2C address configuration (default address \texttt{0x20}).
\end{itemize}

\section{Motors Do Not Move}
\label{sec:motors}

\subsection*{Electrical checks}
\begin{itemize}
    \item Check motor driver power supply and polarity.
    \item Confirm the driver \textbf{logic ground} is connected to ESP32 ground.
    \item Confirm motor outputs are connected to the motor terminals (swap leads if direction is reversed).
\end{itemize}

\subsection*{Signal checks}
\begin{itemize}
    \item Confirm PWM pins match your wiring. In the default pin map, PWM channels include:
    GPIO18/19/23/25/26/27.
    \item Confirm direction pins are connected correctly:
    \texttt{PIN\_MD\_LEFT = GPIO12} and \texttt{PIN\_MD\_RIGHT = GPIO15} (both should be kept low at boot).
    \item If your driver needs two direction pins per motor (IN1/IN2), you must adapt the wiring/firmware accordingly.
\end{itemize}

\subsection*{Practical debugging tip}
Test motor outputs with a simple ``LED + resistor'' on the PWM pin first (you should see brightness change with stick motion),
then connect the motor driver.

\section{Servo Does Not Respond or Jitters}
\label{sec:servo}

\begin{itemize}
    \item Ensure the servo has a proper 5V supply (do \textbf{not} power medium/large servos from the ESP32 board).
    \item Confirm servo ground is tied to ESP32 ground.
    \item Verify the control signal is wired to a configured PWM output pin (e.g., GPIO23 or GPIO25).
    \item If the servo jitters, add a better power supply and reduce electrical noise (shorter wires, capacitors, separate power rails).
\end{itemize}

\section{Telemetry Data Not Appearing in the App}
\label{sec:telemetry}

\subsection*{Symptom: Connected, but no telemetry widgets update}
\begin{itemize}
    \item Wait a few seconds after connecting: the ESP32 sends configuration to the app on first connect.
    \item Verify the sensor is wired to the correct input type:
    \begin{itemize}
        \item Analog telemetry: use ADC1 pins (e.g., GPIO34/35/36/39 in \texttt{FULL\_RC\_MODE}).
        \item Digital indicators: GPIO4/5/2/0 (note GPIO0 boot caution).
        \item I2C sensors: SDA GPIO21, SCL GPIO22.
    \end{itemize}
    \item If you are using \texttt{FULL\_RC\_MODE\_MCP}, confirm the MCP23017 is detected and wired correctly.
\end{itemize}

\subsection*{Symptom: Telemetry is present but values look wrong}
\begin{itemize}
    \item Check sensor scaling (voltage dividers, ADC range, units).
    \item Reduce noise: shorten analog wires, add filtering (RC filter), and avoid running analog lines alongside motor wires.
\end{itemize}

\section{I2C Devices Not Detected (MCP23017 or Sensors)}
\label{sec:i2c}

\begin{itemize}
    \item Verify SDA is on GPIO21 and SCL on GPIO22 (per \texttt{pins.h}).
    \item Confirm the device uses 3.3V logic (or use level shifting).
    \item Check pull-up resistors on SDA/SCL (many breakout boards include them; too many pull-ups in parallel can also cause issues).
    \item Confirm the I2C address (MCP23017 default in the repo is \texttt{0x20}).
\end{itemize}

\section{General Debugging Tips (Recommended Workflow)}
\label{sec:debugging-tips}

\begin{enumerate}
    \item \textbf{Start with the link:} verify Bluetooth pairing and a stable connection first.
    \item \textbf{Test outputs without motors:} LEDs on PWM/switch pins confirm the control path is working.
    \item \textbf{Add one subsystem at a time:} motors, then servos, then sensors, then I2C expansion.
    \item \textbf{Use logging:}
    \begin{itemize}
        \item ESP32: serial debug output (when enabled in the firmware).
        \item Android: Check the indicators to confirm connection state and incoming telemetry.
    \end{itemize}
    \item \textbf{Treat power as a first-class component:} most intermittent failures come from undervoltage or noisy supplies.
\end{enumerate}